\documentclass[11pt]{article}

\usepackage[margin=1in]{geometry}
\usepackage{amsmath,amssymb,amsthm}
\usepackage[hidelinks]{hyperref}

\theoremstyle{plain}
\newtheorem{theorem}{Theorem}
\newtheorem{lemma}[theorem]{Lemma}
\newtheorem{proposition}[theorem]{Proposition}
\newtheorem{corollary}[theorem]{Corollary}
\theoremstyle{definition}
\newtheorem{definition}[theorem]{Definition}
\theoremstyle{remark}
\newtheorem{remark}[theorem]{Remark}

\newcommand{\conjid}{00000002486}
\newcommand{\Par}{\operatorname{Par}}
\newcommand{\dom}{\trianglelefteq}

\title{The dominance order on partitions of $n$ is a self-dual lattice\\[0.4em]\large The Last Math Competition, Conjecture \conjid}
\author{Feiyu\\ \small Independent researcher}
\date{2026-10-02}

\begin{document}
\maketitle

\begin{abstract}
For every $n\ge0$, the partitions of $n$ ordered by dominance form a lattice:
any two partitions have a meet and a join. Conjugation $\lambda\mapsto\lambda'$ is
an involution that reverses dominance, so it exchanges meets and joins and the
lattice is self-dual. This is a classical theorem of Brylawski. We give a short
self-contained proof and a complete Lean~4 formalization on Mathlib's
\texttt{Nat.Partition}.
\end{abstract}

\section{The conjecture as stated}

The original text of conjecture \conjid{} reads:

\begin{quote}
Definition: The order on partitions: dominance. Conjecture: The dominance partial order on partitions is a lattice structure: joins and meets exist, and the structure closes under conjugation duality. (partition order dominance lattice structure)
\end{quote}

\section{Reading of the statement}

% If the statement is ambiguous, under-specified, or contains an error, say so
% HERE and fix a precise reading. State explicitly which reading is being
% settled and why it is the faithful one. Do not silently repair the statement.

A partition $\lambda$ of $n$ is written with its parts in weakly decreasing
order, $\lambda_1\ge\lambda_2\ge\dots\ge\lambda_\ell>0$, and we set $\lambda_i=0$
for $i>\ell$. Its partial sums are
\[
  \Sigma_k(\lambda)=\lambda_1+\dots+\lambda_k\qquad(k\ge0).
\]
The dominance order on the set $\Par(n)$ of partitions of $n$ is
$\lambda\dom\mu$ iff $\Sigma_k(\lambda)\le\Sigma_k(\mu)$ for all $k$. The
conjugate partition $\lambda'$ has parts $\lambda'_j=\#\{i:\lambda_i\ge j\}$
($1\le j\le\lambda_1$), the column lengths of the Young diagram.

Dominance compares partitions of the same integer $n$; this is the standard
setting. So the statement is about the poset $(\Par(n),\dom)$ for each $n$. We
read its three clauses as follows.
\begin{enumerate}
  \item \emph{Partial order}: $\dom$ is reflexive, transitive and
    antisymmetric on $\Par(n)$.
  \item \emph{Lattice}: any $\lambda,\mu\in\Par(n)$ have a greatest lower bound
    (meet) $\lambda\wedge\mu$ and a least upper bound (join) $\lambda\vee\mu$.
  \item \emph{Conjugation duality}: conjugation maps $\Par(n)$ to itself, is an
    involution ($\lambda''=\lambda$) and reverses the order
    ($\lambda\dom\mu\iff\mu'\dom\lambda'$). Hence it is an isomorphism from the
    lattice onto its dual. It exchanges meets and joins:
    $(\lambda\wedge\mu)'=\lambda'\vee\mu'$ and
    $(\lambda\vee\mu)'=\lambda'\wedge\mu'$.
\end{enumerate}

\section{Result}

% State exactly what is established: proved, disproved, or partially resolved.

The conjecture is true. All three clauses hold for every $n\ge0$.

\begin{theorem}\label{thm:main}
Let $n\ge0$.
\begin{enumerate}
  \item[(a)] $(\Par(n),\dom)$ is a partial order.
  \item[(b)] Any $\lambda,\mu\in\Par(n)$ have a meet. It is the partition $\nu$
    with $\Sigma_k(\nu)=\min(\Sigma_k(\lambda),\Sigma_k(\mu))$ for all $k$.
  \item[(c)] $\lambda''=\lambda$, and $\lambda\dom\mu\iff\mu'\dom\lambda'$.
  \item[(d)] $m$ is a meet of $\lambda,\mu$ iff $m'$ is a join of
    $\lambda',\mu'$. In particular any $\lambda,\mu$ have the join
    $\lambda\vee\mu=(\lambda'\wedge\mu')'$, and $(\Par(n),\dom)$ is a lattice
    isomorphic to its dual via conjugation.
\end{enumerate}
\end{theorem}

\section{Proof}

Throughout, $\lambda\in\Par(n)$. Since every part is at least~$1$, there are at
most $n$ parts, so $\Sigma_0(\lambda)=0$ and $\Sigma_k(\lambda)=n$ for
$k\ge n$. The increments $\Sigma_{k+1}(\lambda)-\Sigma_k(\lambda)=\lambda_{k+1}$
are weakly decreasing, so
\begin{equation}\label{eq:concave}
  \Sigma_k(\lambda)+\Sigma_{k+2}(\lambda)\le2\,\Sigma_{k+1}(\lambda)\qquad(k\ge0).
\end{equation}

\begin{lemma}[Realization]\label{lem:real}
Let $S:\mathbb N\to\mathbb N$ be nondecreasing with $S(0)=0$,
$S(k)+S(k+2)\le2S(k+1)$ for all $k$, and $S(k)=n$ for $k\ge n$. Then there is
exactly one $\lambda\in\Par(n)$ with $\Sigma_k(\lambda)=S(k)$ for all $k$.
\end{lemma}

\begin{proof}
Put $d_k=S(k+1)-S(k)\ge0$. The hypothesis says $d_{k+1}\le d_k$, and
$d_n=0$. Let $J$ be the least index with $d_J=0$. Then
$d_0\ge\dots\ge d_{J-1}>0$ and $d_k=0$ for $k\ge J$. So
$d_0+\dots+d_{J-1}=S(n)-S(0)=n$, and the partition $\lambda$ with parts
$d_0,\dots,d_{J-1}$ satisfies $\Sigma_k(\lambda)=d_0+\dots+d_{k-1}=S(k)$.
For uniqueness, the parts of any partition are recovered from its partial sums
as $\lambda_{k+1}=\Sigma_{k+1}(\lambda)-\Sigma_k(\lambda)$.
\end{proof}

\paragraph{(a)} Reflexivity and transitivity are inherited from $\le$ on
$\mathbb N$. If $\lambda\dom\mu\dom\lambda$, then $\Sigma(\lambda)=\Sigma(\mu)$,
so $\lambda=\mu$ by the uniqueness in Lemma~\ref{lem:real}.

\paragraph{(b)} Let $S(k)=\min(\Sigma_k(\lambda),\Sigma_k(\mu))$. It is
nondecreasing, $S(0)=0$, and $S(k)=n$ for $k\ge n$. For concavity, say
$S(k+1)=\Sigma_{k+1}(\lambda)$. Then by \eqref{eq:concave},
\[
  S(k)+S(k+2)\le\Sigma_k(\lambda)+\Sigma_{k+2}(\lambda)\le2\Sigma_{k+1}(\lambda)
  =2S(k+1).
\]
Lemma~\ref{lem:real} gives $\nu$ with $\Sigma(\nu)=S$. Then $\nu\dom\lambda$ and
$\nu\dom\mu$. If $\rho\dom\lambda$ and $\rho\dom\mu$, then
$\Sigma_k(\rho)\le S(k)=\Sigma_k(\nu)$, so $\rho\dom\nu$.

\paragraph{Conjugation and partial sums.} Count the cells of the Young diagram
in the first $t$ columns row by row. This gives
\begin{equation}\label{eq:conjsum}
  \Sigma_t(\lambda')=\sum_i\min(\lambda_i,t)=n-G_\lambda(t),\qquad
  G_\lambda(t):=\sum_i(\lambda_i-t)^+ .
\end{equation}
The second equality uses $\min(x,t)+(x-t)^+=x$.

\begin{lemma}\label{lem:legendre}
For all $J,t\ge0$ we have $\Sigma_J(\lambda)\le Jt+G_\lambda(t)$. Equality holds
whenever $\lambda_i\ge t$ for $i\le J$ and $\lambda_i\le t$ for $i>J$. In
particular:
\begin{enumerate}
  \item[(i)] for every $t$ there is a $J$ with equality, e.g.\
    $J=\#\{i:\lambda_i>t\}$;
  \item[(ii)] for every $J$ there is a $t$ with equality, namely
    $t=\lambda_{J+1}$.
\end{enumerate}
\end{lemma}

\begin{proof}
We have
$\Sigma_J(\lambda)=\sum_{i\le J}\lambda_i\le\sum_{i\le J}\bigl(t+(\lambda_i-t)^+\bigr)
\le Jt+G_\lambda(t)$. Under the stated condition both inequalities are
equalities. In~(i) and~(ii) the condition holds because the parts are weakly
decreasing.
\end{proof}

\paragraph{(c), order reversal.} Let $\lambda\dom\mu$ and $t\ge0$. Choose $J$ as
in Lemma~\ref{lem:legendre}(i) for $\lambda$. Then
\[
  G_\lambda(t)=\Sigma_J(\lambda)-Jt\le\Sigma_J(\mu)-Jt\le G_\mu(t),
\]
and the last step is Lemma~\ref{lem:legendre} for $\mu$. By
\eqref{eq:conjsum}, $\Sigma_t(\mu')\le\Sigma_t(\lambda')$ for every $t$, that is,
$\mu'\dom\lambda'$.

\paragraph{(c), involution.} Fix $k$. By \eqref{eq:conjsum} applied to
$\lambda'$, $\Sigma_k(\lambda'')=n-G_{\lambda'}(k)$, so it suffices to show
$\Sigma_k(\lambda)+G_{\lambda'}(k)=n$.
\begin{itemize}
  \item By Lemma~\ref{lem:legendre}(i) for $\lambda'$ there is $J$ with
    $\Sigma_J(\lambda')=Jk+G_{\lambda'}(k)$, so
    $n-G_\lambda(J)=Jk+G_{\lambda'}(k)$ by \eqref{eq:conjsum}.
    Lemma~\ref{lem:legendre} for $\lambda$ gives
    $\Sigma_k(\lambda)\le kJ+G_\lambda(J)=n-G_{\lambda'}(k)$.
  \item By Lemma~\ref{lem:legendre}(ii) for $\lambda$ there is $t$ with
    $\Sigma_k(\lambda)=kt+G_\lambda(t)$. Lemma~\ref{lem:legendre} for $\lambda'$
    gives $n-G_\lambda(t)=\Sigma_t(\lambda')\le tk+G_{\lambda'}(k)$, so
    $\Sigma_k(\lambda)\ge n-G_{\lambda'}(k)$.
\end{itemize}
Hence $\Sigma(\lambda'')=\Sigma(\lambda)$, and $\lambda''=\lambda$ by~(a).
Applying the order reversal to $\mu'\dom\lambda'$ gives
$\lambda''\dom\mu''$, i.e.\ $\lambda\dom\mu$. This is the converse.

\paragraph{(d)} By~(c), conjugation is a bijection of $\Par(n)$ onto itself
that reverses $\dom$ in both directions. Such a map sends greatest lower
bounds to least upper bounds and conversely. Applying this to the meet
$\lambda'\wedge\mu'$ from~(b) shows that $(\lambda'\wedge\mu')'$ is a join of
$\lambda''=\lambda$ and $\mu''=\mu$. \qed

\section{Formalization}

The accompanying Lean~4 project (\texttt{lean/Submission/Basic.lean}) states
and proves Theorem~\ref{thm:main} for every $n$. Partitions are Mathlib's
\texttt{Nat.Partition n}.

\begin{itemize}
  \item \texttt{sortedParts}, \texttt{partialSum}, \texttt{Dominates} --- the
    parts in weakly decreasing order, the partial sums $\Sigma_k$, and $\dom$.
    \texttt{IsMeet}, \texttt{IsJoin} --- greatest lower and least upper bounds
    for \texttt{Dominates}.
  \item \texttt{colLen p j} $=\#\{i:\lambda_i>j\}$ and \texttt{conj} --- the
    conjugate, whose parts are the nonzero values among
    \texttt{colLen p j}, $0\le j<n$.
  \item \texttt{exists\_of\_concave}, \texttt{eq\_of\_partialSum\_eq} ---
    Lemma~\ref{lem:real}; \texttt{exists\_meet} --- part~(b).
  \item \texttt{partialSum\_conj} --- formula \eqref{eq:conjsum}.
  \item \texttt{partialSum\_le}, \texttt{partialSum\_split} ---
    Lemma~\ref{lem:legendre}.
  \item \texttt{exists\_split\_threshold}, \texttt{exists\_split\_index}
    --- Lemma~\ref{lem:legendre}(i) and~(ii).
  \item \texttt{dominates\_conj}, \texttt{conj\_conj},
    \texttt{dominates\_iff\_conj} --- part~(c);
    \texttt{isMeet\_iff\_isJoin\_conj}, \texttt{exists\_join} --- part~(d).
  \item \texttt{ConjectureHolds} --- for every $n$: reflexivity, transitivity
    and antisymmetry of \texttt{Dominates}; existence of meets and joins;
    \texttt{conj (conj p) = p}; order reversal; and meets of $p,q$
    correspond to joins of their conjugates.
    \texttt{conjecture\_00000002486} proves it.
  \item \texttt{dominanceLattice n} packages the result as a Mathlib
    \texttt{Lattice (Nat.Partition n)} whose order is \texttt{Dominates}.
    \texttt{conj\_sup} and \texttt{conj\_inf} state
    $(\lambda\vee\mu)'=\lambda'\wedge\mu'$ and
    $(\lambda\wedge\mu)'=\lambda'\vee\mu'$.
\end{itemize}

The toolchain is \texttt{leanprover/lean4:v4.33.1}, Mathlib is pinned to
\texttt{v4.33.1}, and \texttt{lake build} succeeds. The project contains no
\texttt{sorry}, no \texttt{admit} and no \texttt{native\_decide}, and
introduces no axioms:
\begin{center}
\texttt{\#print axioms conjecture\_00000002486}
\end{center}
reports exactly \texttt{propext}, \texttt{Classical.choice} and
\texttt{Quot.sound}.

\section{Remarks}

% Optional: what the truth is likely to be, what stays open, why the original
% statement went wrong, what a repaired conjecture should say.

The theorem is due to Brylawski~\cite{bry}. The order reversal under
conjugation is also in Macdonald's book~\cite[Ch.~I, (1.11)]{mac}. The
argument above organizes both through the Legendre-type identity of
Lemma~\ref{lem:legendre}:
$\Sigma_J(\lambda)=\min_t\bigl(Jt+G_\lambda(t)\bigr)$ and
$G_\lambda(t)=\max_J\bigl(\Sigma_J(\lambda)-Jt\bigr)$.

The join is not the pointwise maximum of partial sums. Take
$\lambda=(3,1,1,1)$ and $\mu=(2,2,2)$ in $\Par(6)$. Their partial sums are
$(3,4,5,6)$ and $(2,4,6,6)$. The pointwise maximum $(3,4,6,6)$ has increments
$3,1,2$, which are not decreasing, so it is not the partial-sum sequence of any
partition. The join is $(3,2,1)$, which equals $(\lambda'\wedge\mu')'$ with
$\lambda'=(4,1,1)$ and $\mu'=(3,3)$. The meet is $(2,2,1,1)$, given by the
pointwise minimum $(2,4,5,6)$.

\begin{thebibliography}{9}
\bibitem{bry} T.~Brylawski, \emph{The lattice of integer partitions},
  Discrete Math. 6 (1973), 201--219.
\bibitem{mac} I.~G.~Macdonald, \emph{Symmetric Functions and Hall
  Polynomials}, 2nd ed., Oxford University Press, 1995.
\end{thebibliography}

\end{document}
